% !TEX root = master.tex
\chapter{Funktionale Umsetzung: Funktion und Skript}
\label{ch:b-funktionen}

Analog zu Projekt~A ordnet dieses Kapitel die Funktionalit\"aten den jeweiligen Skripten zu
(Tabelle~\ref{tab:b-funktionen}).

\begingroup
\small
\setlength{\tabcolsep}{5pt}
\renewcommand{\arraystretch}{1.3}
\begin{longtable}{@{}L{4.4cm} L{4.4cm} L{4.7cm}@{}}
\caption{Zuordnung Funktionalit\"at zu Skript (Projekt~B)}
\label{tab:b-funktionen}\\
\toprule
\textbf{Funktionalit\"at} & \textbf{Skript / Datei} & \textbf{Kurzbeschreibung} \\
\midrule
\endfirsthead
\multicolumn{3}{@{}l}{\small\itshape Tabelle \thetable{} -- Fortsetzung}\\
\midrule
\textbf{Funktionalit\"at} & \textbf{Skript / Datei} & \textbf{Kurzbeschreibung} \\
\midrule
\endhead
\midrule
\multicolumn{3}{r@{}}{\footnotesize\itshape Fortsetzung n\"achste Seite}\\
\endfoot
\bottomrule
\endlastfoot

\multicolumn{3}{@{}l}{\textbf{Buchung}}\\*
\addlinespace[0.4ex]
Buchungsformular & \texttt{booking.php} & Eingabemaske f\"ur die Buchung \\
Echtzeit-Kostenberechnung & \texttt{js/booking.js} & Personen $\times$ Tage $\times$ Preis + Ausfl\"uge \\
Dynamische Personenfelder & \texttt{js/booking.js} & Felder je nach Personenzahl \\
Clientseitige Validierung & \texttt{js/booking.js} & Pflichtfelder und Datumspr\"ufung \\
Buchungsverarbeitung & \texttt{process\_booking.php} & POST auswerten, Preis berechnen, Kartennummer maskieren \\
Datenspeicherung (JSON + Zeitstempel) & \texttt{process\_booking.php} & Ablage je Buchung unter \texttt{bookings/} \\
Buchungsbest\"atigung & \texttt{confirmation.php} & Zusammenfassung aus der Session \\
GET-/POST-Demonstration & \texttt{formcheck.php} & Test der Formular\"ubertragung \\
\addlinespace[0.8ex]

\multicolumn{3}{@{}l}{\textbf{Oberfl\"ache und Navigation}}\\*
\addlinespace[0.4ex]
Hell-/Dunkel-Modus & \texttt{js/main.js} & Umschaltung, gespeichert im \texttt{localStorage} \\
Mobiles Men\"u / Navigation & \texttt{js/main.js}, \texttt{includes/header.php} & responsive Navigation \\
News-Ticker / Besucherz\"ahler & \texttt{includes/footer.php} & Z\"ahlung \"uber \texttt{visitor\_count.txt} \\
Gemeinsamer Kopf-/Fu\ss{}bereich & \texttt{includes/header.php}, \texttt{includes/footer.php} & wiederverwendbare Bausteine \\
Gestaltung / Theme & \texttt{css/style.css} & Farbschema, Layout, Hell-/Dunkel-Modus \\
\addlinespace[0.8ex]

\multicolumn{3}{@{}l}{\textbf{Inhaltsseiten}}\\*
\addlinespace[0.4ex]
Kundenstimmen & \texttt{tours.php} & mehrere Bewertungen \\
Teamseite & \texttt{about.php} & Vorstellung des fiktiven Teams \\
Rechtliche Seiten & \texttt{legal.php} & Impressum, Datenschutz, AGB \\
\end{longtable}
\endgroup

\section{Echtzeit-Kostenberechnung}
Die Buchungsseite berechnet den Gesamtpreis ohne Neuladen der Seite. Bei jeder \"Anderung der
Eingaben (Personenzahl, Reisedauer, ausgew\"ahlte Ausfl\"uge) wird der Preis neu ermittelt und als
Zusammenfassung angezeigt. Die Grundkosten ergeben sich aus Personenzahl, Anzahl der Tage und dem
Tagespreis; hinzu kommen die Kosten der gew\"ahlten Ausfl\"uge.

\begin{lstlisting}[language=Java,caption={Kern der Kostenberechnung (Auszug \texttt{js/booking.js})},label={lst:b-cost}]
var baseCost = persons * days * pricePerDay;     // Grundpreis
var excursionCost = 0;                           // Summe der Ausfluege
// ... Aufsummieren der ausgewaehlten Ausfluege ...
var total = baseCost + excursionCost;            // Gesamtpreis
\end{lstlisting}

\section{Speicherung der Buchungen}
Bei Absenden des Formulars wertet \texttt{process\_booking.php} die \"ubermittelten Daten aus,
berechnet den Gesamtpreis erneut serverseitig und maskiert die angegebene Kreditkartennummer.
Jede Buchung wird anschlie\ss{}end mit einem Zeitstempel als eigene \ac{json}-Datei im Verzeichnis
\texttt{bookings/} gespeichert (Dateiname nach dem Muster \texttt{ST-JJJJMMTT-XXXXXX.json}). Die
Zusammenfassung wird \"uber die Session an die Best\"atigungsseite \"ubergeben.
